\cabstract{
在学习王老师的大数据原理与技术课程\cite{lecture01}后，在此报告中完成相应的课后作业，本次作业（作业-6）任务为图文多模态问答（VQA），要求使用 ViLT 预训练模型（通过 Hugging Face 库）进行调用与测试，可进行微调或直接推理。

实验要求：
\begin{itemize}
    \item 使用 ViLT 预训练模型（Hugging Face）。
    \item 尝试调用模型，输入图片与问题（例如“图中有什么动物？”），也可自定义问题，并输出模型答案。
\end{itemize}

作业提交及注意事项：
\begin{itemize}
    \item 提交内容：代码与简要报告（报告需包含测试案例截图，至少 2 页）。
    \item 命名格式：学号+姓名+作业6.zip
    \item 提交邮箱：BDSysu2026Spring@163.com
    \item 截止时间：布置作业之后的下一个周日 24:00
    \item 严禁抄袭：任何一次抄袭，平时成绩按零分处理
\end{itemize}

本实验首先使用 Hugging Face 提供的 ViLT 预训练模型及其处理器对输入图像和问题进行前向推理。对 CIFAR-10 子集进行快速验证时，图像先被转换为 PIL.Image 并按模型要求缩放与归一化（如 Resize 到模型输入尺寸），问题使用通用模板（例如“What is in the image?”）或自定义问题。推理得到的自由文本答案被保存为原始输出，随后对结果做小写化、去除标点、同义词替换与基于相似度的模糊匹配，将开放域答案映射到 CIFAR-10 的 10 类标签。作为基线比较，使用预训练的 ResNet-18 对目标类别进行监督微调（测试子集约 1000 张，batch=64，epochs=3，lr=1e-3）。

结论：预训练的 ViLT 在开放域 VQA 上能给出有意义的自由文本答案，但在将自由文本映射到固定标签集合时存在误匹配风险。本实验表明，针对目标类别进行监督微调（例如 ResNet-18）或结合候选重排序与同义词映射的后处理策略，可显著提高分类准确率与一致性。
}
\ckeywords{图文多模态问答，VQA，ViLT，Hugging Face}
\ekeywords{}


